% =====================================================================
% Paper 14 — Conditional BRST/FP closure note for the TECT defect sector
%   (HEALTHIEST of Paper-12/13/14; GAP-2 = T6 already in canonical
%    bookkeeping; over-claim tone reduction only)
% Status: [DRAFT] [NEEDS_UPDATE]
% AUDIT-FLAG: AUDIT-2026-05-02-Wave7-Aux-Epoch-Overclaim
%   (Math314-AddD extension covering Wave 1/4/5 papers Paper-09..16).
%   Hostile-referee audit (2026-05-02): canonical alignment OK
%   (GAP-2 = T6 in Math307 / Math310 bookkeeping). Three light fixes:
%   (1) FP determinant / ghost path-integral formula sketch needs
%       textbook-standard Grassmann measure + sign convention; current
%       sketch too short.
%   (2) Berry-correction inheritance from Math160 not fully derived;
%       paper inherits result rather than re-deriving.
%   (3) "GAP-2 is now complete" / "fully rigorous" / "ready for
%       publication" wording too strong. Recast as conditional BRST/FP
%       closure note.
% Compile: pdflatex Paper-14.tex (×2 for refs)
% Canonical archive: Math160, Math280, Math307, Math310, Math314-AddD
% =====================================================================
% --- TECT canonical preamble (see _shared/tect-preamble.tex) ----------
% =====================================================================
% TECT canonical LaTeX preamble (revtex4-2 / single-column / theorem-safe)
% =====================================================================
% Maintainer: Jusang Lee (jtkor@outlook.com)
% Binding from: 2026-05-06
% Purpose: a single source-of-truth preamble for every TECT paper
% (Paper-00 .. Paper-10 + Auxiliary notes). Designed to (a) eliminate
% font-corruption on pdflatex (T1 fontenc + lmodern), (b) make
% theorem-heavy mathematical-physics content render cleanly in a
% single-column layout (PRD class option, NOT PRL two-column), (c)
% provide consistent theorem numbering across all TECT papers via
% amsthm with a shared counter set, (d) make hyperref + cleveref
% cross-references resolve cleanly with the standard two-pass
% pdflatex compile.
%
% Usage in each Paper-NN.tex:
%   % =====================================================================
% TECT canonical LaTeX preamble (revtex4-2 / single-column / theorem-safe)
% =====================================================================
% Maintainer: Jusang Lee (jtkor@outlook.com)
% Binding from: 2026-05-06
% Purpose: a single source-of-truth preamble for every TECT paper
% (Paper-00 .. Paper-10 + Auxiliary notes). Designed to (a) eliminate
% font-corruption on pdflatex (T1 fontenc + lmodern), (b) make
% theorem-heavy mathematical-physics content render cleanly in a
% single-column layout (PRD class option, NOT PRL two-column), (c)
% provide consistent theorem numbering across all TECT papers via
% amsthm with a shared counter set, (d) make hyperref + cleveref
% cross-references resolve cleanly with the standard two-pass
% pdflatex compile.
%
% Usage in each Paper-NN.tex:
%   % =====================================================================
% TECT canonical LaTeX preamble (revtex4-2 / single-column / theorem-safe)
% =====================================================================
% Maintainer: Jusang Lee (jtkor@outlook.com)
% Binding from: 2026-05-06
% Purpose: a single source-of-truth preamble for every TECT paper
% (Paper-00 .. Paper-10 + Auxiliary notes). Designed to (a) eliminate
% font-corruption on pdflatex (T1 fontenc + lmodern), (b) make
% theorem-heavy mathematical-physics content render cleanly in a
% single-column layout (PRD class option, NOT PRL two-column), (c)
% provide consistent theorem numbering across all TECT papers via
% amsthm with a shared counter set, (d) make hyperref + cleveref
% cross-references resolve cleanly with the standard two-pass
% pdflatex compile.
%
% Usage in each Paper-NN.tex:
%   \input{../_shared/tect-preamble.tex}
%   \begin{document}
%   ...
%   \end{document}
%
% Build (every paper): `pdflatex Paper-NN && pdflatex Paper-NN`
% (the second pass resolves \ref{} / \cref{} forward references).
% A bash/PowerShell wrapper that automates this is at
% Docs/papers/papers/_shared/build-paper.sh / build-paper.ps1.
% =====================================================================

% ---- Document class --------------------------------------------------
% PRD class option of revtex4-2 with [reprint,onecolumn] gives the same
% APS heading / by-line / abstract style as PRL but in a wider single
% column that handles long display equations and theorem environments
% without column-break artefacts. nofootinbib keeps citations out of
% footnotes. The 11pt option restores standard font metrics; with T1
% fontenc + lmodern this gives non-bitmap glyphs.
\documentclass[%
  aps,prd,reprint,onecolumn,nofootinbib,11pt,%
  amsmath,amssymb,floatfix,longbibliography%
]{revtex4-2}

% ---- Encoding & fonts ------------------------------------------------
\usepackage[T1]{fontenc}        % proper 8-bit font encoding (no font corruption)
\usepackage[utf8]{inputenc}     % allow UTF-8 source
\usepackage{lmodern}            % scalable Latin Modern (replaces bitmap CM)
\usepackage{textcomp}           % extra text symbols
\usepackage{microtype}          % subtle kerning improvements

% ---- UTF-8 → LaTeX character mappings --------------------------------
% Maps the Unicode characters that occur in TECT body text (legacy from
% the chat-content auto-archival rule, CLAUDE.md §4) to LaTeX-safe
% commands. Without these declarations, T1+inputenc rejects the chars
% with "Unicode character X not set up for use with LaTeX".
% Adding declarations centrally (here) is preferable to per-file edits.
\DeclareUnicodeCharacter{00A7}{\S}              % §  (section sign)
\DeclareUnicodeCharacter{00B2}{\ensuremath{^{2}}}        % ²
\DeclareUnicodeCharacter{00B3}{\ensuremath{^{3}}}        % ³
\DeclareUnicodeCharacter{00B5}{\ensuremath{\mu}}         % µ (micro)
\DeclareUnicodeCharacter{00B7}{\ensuremath{\cdot}}       % · (middle dot)
\DeclareUnicodeCharacter{00D7}{\ensuremath{\times}}      % ×
\DeclareUnicodeCharacter{00E4}{\"a}             % ä
\DeclareUnicodeCharacter{00E9}{\'e}             % é
\DeclareUnicodeCharacter{00F6}{\"o}             % ö
\DeclareUnicodeCharacter{00FC}{\"u}             % ü
\DeclareUnicodeCharacter{010C}{\v{C}}           % Č
\DeclareUnicodeCharacter{010D}{\v{c}}           % č
\DeclareUnicodeCharacter{0394}{\ensuremath{\Delta}}      % Δ
\DeclareUnicodeCharacter{03A3}{\ensuremath{\Sigma}}      % Σ
\DeclareUnicodeCharacter{03A9}{\ensuremath{\Omega}}      % Ω
\DeclareUnicodeCharacter{03B1}{\ensuremath{\alpha}}      % α
\DeclareUnicodeCharacter{03B2}{\ensuremath{\beta}}       % β
\DeclareUnicodeCharacter{03B3}{\ensuremath{\gamma}}      % γ
\DeclareUnicodeCharacter{03B4}{\ensuremath{\delta}}      % δ
\DeclareUnicodeCharacter{03B5}{\ensuremath{\varepsilon}} % ε
\DeclareUnicodeCharacter{03BA}{\ensuremath{\kappa}}      % κ
\DeclareUnicodeCharacter{03BB}{\ensuremath{\lambda}}     % λ
\DeclareUnicodeCharacter{03BC}{\ensuremath{\mu}}         % μ
\DeclareUnicodeCharacter{03BD}{\ensuremath{\nu}}         % ν
\DeclareUnicodeCharacter{03C0}{\ensuremath{\pi}}         % π
\DeclareUnicodeCharacter{03C1}{\ensuremath{\rho}}        % ρ
\DeclareUnicodeCharacter{03C3}{\ensuremath{\sigma}}      % σ
\DeclareUnicodeCharacter{03C4}{\ensuremath{\tau}}        % τ
\DeclareUnicodeCharacter{03C6}{\ensuremath{\varphi}}     % φ
\DeclareUnicodeCharacter{03C7}{\ensuremath{\chi}}        % χ
\DeclareUnicodeCharacter{03C8}{\ensuremath{\psi}}        % ψ
\DeclareUnicodeCharacter{03C9}{\ensuremath{\omega}}      % ω
\DeclareUnicodeCharacter{2013}{--}              % – (en dash)
\DeclareUnicodeCharacter{2014}{---}             % — (em dash)
\DeclareUnicodeCharacter{2018}{`}               % ‘
\DeclareUnicodeCharacter{2019}{'}               % ’
\DeclareUnicodeCharacter{201C}{``}              % “
\DeclareUnicodeCharacter{201D}{''}              % ”
\DeclareUnicodeCharacter{2070}{\ensuremath{^{0}}}        % ⁰
\DeclareUnicodeCharacter{2071}{\ensuremath{^{i}}}        % ⁱ
\DeclareUnicodeCharacter{2122}{\textsuperscript{TM}}     % ™
\DeclareUnicodeCharacter{2126}{\ensuremath{\Omega}}      % Ω (ohm sign)
\DeclareUnicodeCharacter{210F}{\ensuremath{\hbar}}       % ℏ
\DeclareUnicodeCharacter{2115}{\ensuremath{\mathbb{N}}}  % ℕ
\DeclareUnicodeCharacter{2102}{\ensuremath{\mathbb{C}}}  % ℂ
\DeclareUnicodeCharacter{211D}{\ensuremath{\mathbb{R}}}  % ℝ
\DeclareUnicodeCharacter{2124}{\ensuremath{\mathbb{Z}}}  % ℤ
\DeclareUnicodeCharacter{211A}{\ensuremath{\mathbb{Q}}}  % ℚ
\DeclareUnicodeCharacter{2192}{\ensuremath{\to}}         % →
\DeclareUnicodeCharacter{21D2}{\ensuremath{\Rightarrow}} % ⇒
\DeclareUnicodeCharacter{2200}{\ensuremath{\forall}}     % ∀
\DeclareUnicodeCharacter{2203}{\ensuremath{\exists}}     % ∃
\DeclareUnicodeCharacter{2208}{\ensuremath{\in}}         % ∈
\DeclareUnicodeCharacter{2209}{\ensuremath{\notin}}      % ∉
\DeclareUnicodeCharacter{2227}{\ensuremath{\wedge}}      % ∧
\DeclareUnicodeCharacter{2228}{\ensuremath{\vee}}        % ∨
\DeclareUnicodeCharacter{2229}{\ensuremath{\cap}}        % ∩
\DeclareUnicodeCharacter{222A}{\ensuremath{\cup}}        % ∪
\DeclareUnicodeCharacter{2245}{\ensuremath{\cong}}       % ≅
\DeclareUnicodeCharacter{2248}{\ensuremath{\approx}}     % ≈
\DeclareUnicodeCharacter{2260}{\ensuremath{\neq}}        % ≠
\DeclareUnicodeCharacter{2264}{\ensuremath{\leq}}        % ≤
\DeclareUnicodeCharacter{2265}{\ensuremath{\geq}}        % ≥
\DeclareUnicodeCharacter{22C5}{\ensuremath{\cdot}}       % ⋅
\DeclareUnicodeCharacter{2295}{\ensuremath{\oplus}}      % ⊕
\DeclareUnicodeCharacter{2297}{\ensuremath{\otimes}}     % ⊗

% ---- Mathematics & symbols ------------------------------------------
\usepackage{amsmath,amssymb,bm,mathrsfs,mathtools}
\usepackage{amsthm}             % theorem environments (load BEFORE hyperref)

% ---- Theorem environments (shared counter set, TECT canonical) -------
% All theorems / lemmas / propositions / corollaries / remarks share a
% single counter so cross-references read consistently across a paper.
% Definitions get a separate counter (definitions are numbered in their
% own sequence by editorial convention).
\newtheorem{theorem}{Theorem}
\newtheorem{lemma}[theorem]{Lemma}
\newtheorem{proposition}[theorem]{Proposition}
\newtheorem{corollary}[theorem]{Corollary}

\theoremstyle{remark}
\newtheorem{remark}[theorem]{Remark}
\newtheorem{example}[theorem]{Example}

\theoremstyle{definition}
\newtheorem{definition}[theorem]{Definition}
\newtheorem{conjecture}[theorem]{Conjecture}
\newtheorem{hypothesis}[theorem]{Hypothesis}

% ---- Tables / floats / graphics --------------------------------------
\usepackage{booktabs}           % \toprule / \midrule / \bottomrule
\usepackage{array}
\usepackage{graphicx}
\usepackage{enumitem}           % \begin{itemize}[leftmargin=...] options
\usepackage{hyphenat}           % \hyp{} for breakable hyphens in \texttt
\usepackage{seqsplit}           % \seqsplit{} for long unbreakable strings
\sloppy                         % allow stretchier inter-word spacing to
                                % reduce overfull \hbox warnings on long URLs
                                % and long \texttt tokens (paragraph-level).
\emergencystretch=3em           % last-resort hbox stretch budget

% ---- Cross-references (hyperref last among the standard packages,
%      cleveref AFTER hyperref) ----------------------------------------
\usepackage[%
  colorlinks=true,%
  linkcolor=blue,%
  citecolor=blue,%
  urlcolor=blue,%
  breaklinks=true,%
  bookmarks=true,%
  bookmarksnumbered=true%
]{hyperref}
\usepackage[capitalise,nameinlink]{cleveref}

% ---- TECT-canonical author block macros -----------------------------
% (defined here so every paper has the same author / affiliation style)
\newcommand{\tectauthor}{%
  \author{Jusang Lee}%
  \email{jtkor@outlook.com}%
  \affiliation{Independent researcher, Republic of Korea}%
}

% ---- TECT-canonical archive footer ----------------------------------
\newcommand{\tectarchivefooter}{%
  \par\medskip
  \noindent\small\textit{Canonical archive}:
  \url{https://github.com/TECT-OpenPhysics/TECT}.\par%
}

% ---- TECT TODO / AUDIT-FLAG / HONEST-SCOPE inline marks --------------
% Used in drafts to highlight pending audit items without disturbing
% the body text style. Suppress via \tectsuppressmarks (define empty).
\providecommand{\tectauditflag}[1]{\textcolor{red}{\textbf{[AUDIT-FLAG: #1]}}}
\providecommand{\tecttodo}[1]{\textcolor{orange}{\textbf{[TODO: #1]}}}
\providecommand{\tecthonestscope}[1]{\textit{Honest scope: #1.}}

% =====================================================================
% End of TECT canonical preamble
% =====================================================================

%   \begin{document}
%   ...
%   \end{document}
%
% Build (every paper): `pdflatex Paper-NN && pdflatex Paper-NN`
% (the second pass resolves \ref{} / \cref{} forward references).
% A bash/PowerShell wrapper that automates this is at
% Docs/papers/papers/_shared/build-paper.sh / build-paper.ps1.
% =====================================================================

% ---- Document class --------------------------------------------------
% PRD class option of revtex4-2 with [reprint,onecolumn] gives the same
% APS heading / by-line / abstract style as PRL but in a wider single
% column that handles long display equations and theorem environments
% without column-break artefacts. nofootinbib keeps citations out of
% footnotes. The 11pt option restores standard font metrics; with T1
% fontenc + lmodern this gives non-bitmap glyphs.
\documentclass[%
  aps,prd,reprint,onecolumn,nofootinbib,11pt,%
  amsmath,amssymb,floatfix,longbibliography%
]{revtex4-2}

% ---- Encoding & fonts ------------------------------------------------
\usepackage[T1]{fontenc}        % proper 8-bit font encoding (no font corruption)
\usepackage[utf8]{inputenc}     % allow UTF-8 source
\usepackage{lmodern}            % scalable Latin Modern (replaces bitmap CM)
\usepackage{textcomp}           % extra text symbols
\usepackage{microtype}          % subtle kerning improvements

% ---- UTF-8 → LaTeX character mappings --------------------------------
% Maps the Unicode characters that occur in TECT body text (legacy from
% the chat-content auto-archival rule, CLAUDE.md §4) to LaTeX-safe
% commands. Without these declarations, T1+inputenc rejects the chars
% with "Unicode character X not set up for use with LaTeX".
% Adding declarations centrally (here) is preferable to per-file edits.
\DeclareUnicodeCharacter{00A7}{\S}              % §  (section sign)
\DeclareUnicodeCharacter{00B2}{\ensuremath{^{2}}}        % ²
\DeclareUnicodeCharacter{00B3}{\ensuremath{^{3}}}        % ³
\DeclareUnicodeCharacter{00B5}{\ensuremath{\mu}}         % µ (micro)
\DeclareUnicodeCharacter{00B7}{\ensuremath{\cdot}}       % · (middle dot)
\DeclareUnicodeCharacter{00D7}{\ensuremath{\times}}      % ×
\DeclareUnicodeCharacter{00E4}{\"a}             % ä
\DeclareUnicodeCharacter{00E9}{\'e}             % é
\DeclareUnicodeCharacter{00F6}{\"o}             % ö
\DeclareUnicodeCharacter{00FC}{\"u}             % ü
\DeclareUnicodeCharacter{010C}{\v{C}}           % Č
\DeclareUnicodeCharacter{010D}{\v{c}}           % č
\DeclareUnicodeCharacter{0394}{\ensuremath{\Delta}}      % Δ
\DeclareUnicodeCharacter{03A3}{\ensuremath{\Sigma}}      % Σ
\DeclareUnicodeCharacter{03A9}{\ensuremath{\Omega}}      % Ω
\DeclareUnicodeCharacter{03B1}{\ensuremath{\alpha}}      % α
\DeclareUnicodeCharacter{03B2}{\ensuremath{\beta}}       % β
\DeclareUnicodeCharacter{03B3}{\ensuremath{\gamma}}      % γ
\DeclareUnicodeCharacter{03B4}{\ensuremath{\delta}}      % δ
\DeclareUnicodeCharacter{03B5}{\ensuremath{\varepsilon}} % ε
\DeclareUnicodeCharacter{03BA}{\ensuremath{\kappa}}      % κ
\DeclareUnicodeCharacter{03BB}{\ensuremath{\lambda}}     % λ
\DeclareUnicodeCharacter{03BC}{\ensuremath{\mu}}         % μ
\DeclareUnicodeCharacter{03BD}{\ensuremath{\nu}}         % ν
\DeclareUnicodeCharacter{03C0}{\ensuremath{\pi}}         % π
\DeclareUnicodeCharacter{03C1}{\ensuremath{\rho}}        % ρ
\DeclareUnicodeCharacter{03C3}{\ensuremath{\sigma}}      % σ
\DeclareUnicodeCharacter{03C4}{\ensuremath{\tau}}        % τ
\DeclareUnicodeCharacter{03C6}{\ensuremath{\varphi}}     % φ
\DeclareUnicodeCharacter{03C7}{\ensuremath{\chi}}        % χ
\DeclareUnicodeCharacter{03C8}{\ensuremath{\psi}}        % ψ
\DeclareUnicodeCharacter{03C9}{\ensuremath{\omega}}      % ω
\DeclareUnicodeCharacter{2013}{--}              % – (en dash)
\DeclareUnicodeCharacter{2014}{---}             % — (em dash)
\DeclareUnicodeCharacter{2018}{`}               % ‘
\DeclareUnicodeCharacter{2019}{'}               % ’
\DeclareUnicodeCharacter{201C}{``}              % “
\DeclareUnicodeCharacter{201D}{''}              % ”
\DeclareUnicodeCharacter{2070}{\ensuremath{^{0}}}        % ⁰
\DeclareUnicodeCharacter{2071}{\ensuremath{^{i}}}        % ⁱ
\DeclareUnicodeCharacter{2122}{\textsuperscript{TM}}     % ™
\DeclareUnicodeCharacter{2126}{\ensuremath{\Omega}}      % Ω (ohm sign)
\DeclareUnicodeCharacter{210F}{\ensuremath{\hbar}}       % ℏ
\DeclareUnicodeCharacter{2115}{\ensuremath{\mathbb{N}}}  % ℕ
\DeclareUnicodeCharacter{2102}{\ensuremath{\mathbb{C}}}  % ℂ
\DeclareUnicodeCharacter{211D}{\ensuremath{\mathbb{R}}}  % ℝ
\DeclareUnicodeCharacter{2124}{\ensuremath{\mathbb{Z}}}  % ℤ
\DeclareUnicodeCharacter{211A}{\ensuremath{\mathbb{Q}}}  % ℚ
\DeclareUnicodeCharacter{2192}{\ensuremath{\to}}         % →
\DeclareUnicodeCharacter{21D2}{\ensuremath{\Rightarrow}} % ⇒
\DeclareUnicodeCharacter{2200}{\ensuremath{\forall}}     % ∀
\DeclareUnicodeCharacter{2203}{\ensuremath{\exists}}     % ∃
\DeclareUnicodeCharacter{2208}{\ensuremath{\in}}         % ∈
\DeclareUnicodeCharacter{2209}{\ensuremath{\notin}}      % ∉
\DeclareUnicodeCharacter{2227}{\ensuremath{\wedge}}      % ∧
\DeclareUnicodeCharacter{2228}{\ensuremath{\vee}}        % ∨
\DeclareUnicodeCharacter{2229}{\ensuremath{\cap}}        % ∩
\DeclareUnicodeCharacter{222A}{\ensuremath{\cup}}        % ∪
\DeclareUnicodeCharacter{2245}{\ensuremath{\cong}}       % ≅
\DeclareUnicodeCharacter{2248}{\ensuremath{\approx}}     % ≈
\DeclareUnicodeCharacter{2260}{\ensuremath{\neq}}        % ≠
\DeclareUnicodeCharacter{2264}{\ensuremath{\leq}}        % ≤
\DeclareUnicodeCharacter{2265}{\ensuremath{\geq}}        % ≥
\DeclareUnicodeCharacter{22C5}{\ensuremath{\cdot}}       % ⋅
\DeclareUnicodeCharacter{2295}{\ensuremath{\oplus}}      % ⊕
\DeclareUnicodeCharacter{2297}{\ensuremath{\otimes}}     % ⊗

% ---- Mathematics & symbols ------------------------------------------
\usepackage{amsmath,amssymb,bm,mathrsfs,mathtools}
\usepackage{amsthm}             % theorem environments (load BEFORE hyperref)

% ---- Theorem environments (shared counter set, TECT canonical) -------
% All theorems / lemmas / propositions / corollaries / remarks share a
% single counter so cross-references read consistently across a paper.
% Definitions get a separate counter (definitions are numbered in their
% own sequence by editorial convention).
\newtheorem{theorem}{Theorem}
\newtheorem{lemma}[theorem]{Lemma}
\newtheorem{proposition}[theorem]{Proposition}
\newtheorem{corollary}[theorem]{Corollary}

\theoremstyle{remark}
\newtheorem{remark}[theorem]{Remark}
\newtheorem{example}[theorem]{Example}

\theoremstyle{definition}
\newtheorem{definition}[theorem]{Definition}
\newtheorem{conjecture}[theorem]{Conjecture}
\newtheorem{hypothesis}[theorem]{Hypothesis}

% ---- Tables / floats / graphics --------------------------------------
\usepackage{booktabs}           % \toprule / \midrule / \bottomrule
\usepackage{array}
\usepackage{graphicx}
\usepackage{enumitem}           % \begin{itemize}[leftmargin=...] options
\usepackage{hyphenat}           % \hyp{} for breakable hyphens in \texttt
\usepackage{seqsplit}           % \seqsplit{} for long unbreakable strings
\sloppy                         % allow stretchier inter-word spacing to
                                % reduce overfull \hbox warnings on long URLs
                                % and long \texttt tokens (paragraph-level).
\emergencystretch=3em           % last-resort hbox stretch budget

% ---- Cross-references (hyperref last among the standard packages,
%      cleveref AFTER hyperref) ----------------------------------------
\usepackage[%
  colorlinks=true,%
  linkcolor=blue,%
  citecolor=blue,%
  urlcolor=blue,%
  breaklinks=true,%
  bookmarks=true,%
  bookmarksnumbered=true%
]{hyperref}
\usepackage[capitalise,nameinlink]{cleveref}

% ---- TECT-canonical author block macros -----------------------------
% (defined here so every paper has the same author / affiliation style)
\newcommand{\tectauthor}{%
  \author{Jusang Lee}%
  \email{jtkor@outlook.com}%
  \affiliation{Independent researcher, Republic of Korea}%
}

% ---- TECT-canonical archive footer ----------------------------------
\newcommand{\tectarchivefooter}{%
  \par\medskip
  \noindent\small\textit{Canonical archive}:
  \url{https://github.com/TECT-OpenPhysics/TECT}.\par%
}

% ---- TECT TODO / AUDIT-FLAG / HONEST-SCOPE inline marks --------------
% Used in drafts to highlight pending audit items without disturbing
% the body text style. Suppress via \tectsuppressmarks (define empty).
\providecommand{\tectauditflag}[1]{\textcolor{red}{\textbf{[AUDIT-FLAG: #1]}}}
\providecommand{\tecttodo}[1]{\textcolor{orange}{\textbf{[TODO: #1]}}}
\providecommand{\tecthonestscope}[1]{\textit{Honest scope: #1.}}

% =====================================================================
% End of TECT canonical preamble
% =====================================================================

%   \begin{document}
%   ...
%   \end{document}
%
% Build (every paper): `pdflatex Paper-NN && pdflatex Paper-NN`
% (the second pass resolves \ref{} / \cref{} forward references).
% A bash/PowerShell wrapper that automates this is at
% Docs/papers/papers/_shared/build-paper.sh / build-paper.ps1.
% =====================================================================

% ---- Document class --------------------------------------------------
% PRD class option of revtex4-2 with [reprint,onecolumn] gives the same
% APS heading / by-line / abstract style as PRL but in a wider single
% column that handles long display equations and theorem environments
% without column-break artefacts. nofootinbib keeps citations out of
% footnotes. The 11pt option restores standard font metrics; with T1
% fontenc + lmodern this gives non-bitmap glyphs.
\documentclass[%
  aps,prd,reprint,onecolumn,nofootinbib,11pt,%
  amsmath,amssymb,floatfix,longbibliography%
]{revtex4-2}

% ---- Encoding & fonts ------------------------------------------------
\usepackage[T1]{fontenc}        % proper 8-bit font encoding (no font corruption)
\usepackage[utf8]{inputenc}     % allow UTF-8 source
\usepackage{lmodern}            % scalable Latin Modern (replaces bitmap CM)
\usepackage{textcomp}           % extra text symbols
\usepackage{microtype}          % subtle kerning improvements

% ---- UTF-8 → LaTeX character mappings --------------------------------
% Maps the Unicode characters that occur in TECT body text (legacy from
% the chat-content auto-archival rule, CLAUDE.md §4) to LaTeX-safe
% commands. Without these declarations, T1+inputenc rejects the chars
% with "Unicode character X not set up for use with LaTeX".
% Adding declarations centrally (here) is preferable to per-file edits.
\DeclareUnicodeCharacter{00A7}{\S}              % §  (section sign)
\DeclareUnicodeCharacter{00B2}{\ensuremath{^{2}}}        % ²
\DeclareUnicodeCharacter{00B3}{\ensuremath{^{3}}}        % ³
\DeclareUnicodeCharacter{00B5}{\ensuremath{\mu}}         % µ (micro)
\DeclareUnicodeCharacter{00B7}{\ensuremath{\cdot}}       % · (middle dot)
\DeclareUnicodeCharacter{00D7}{\ensuremath{\times}}      % ×
\DeclareUnicodeCharacter{00E4}{\"a}             % ä
\DeclareUnicodeCharacter{00E9}{\'e}             % é
\DeclareUnicodeCharacter{00F6}{\"o}             % ö
\DeclareUnicodeCharacter{00FC}{\"u}             % ü
\DeclareUnicodeCharacter{010C}{\v{C}}           % Č
\DeclareUnicodeCharacter{010D}{\v{c}}           % č
\DeclareUnicodeCharacter{0394}{\ensuremath{\Delta}}      % Δ
\DeclareUnicodeCharacter{03A3}{\ensuremath{\Sigma}}      % Σ
\DeclareUnicodeCharacter{03A9}{\ensuremath{\Omega}}      % Ω
\DeclareUnicodeCharacter{03B1}{\ensuremath{\alpha}}      % α
\DeclareUnicodeCharacter{03B2}{\ensuremath{\beta}}       % β
\DeclareUnicodeCharacter{03B3}{\ensuremath{\gamma}}      % γ
\DeclareUnicodeCharacter{03B4}{\ensuremath{\delta}}      % δ
\DeclareUnicodeCharacter{03B5}{\ensuremath{\varepsilon}} % ε
\DeclareUnicodeCharacter{03BA}{\ensuremath{\kappa}}      % κ
\DeclareUnicodeCharacter{03BB}{\ensuremath{\lambda}}     % λ
\DeclareUnicodeCharacter{03BC}{\ensuremath{\mu}}         % μ
\DeclareUnicodeCharacter{03BD}{\ensuremath{\nu}}         % ν
\DeclareUnicodeCharacter{03C0}{\ensuremath{\pi}}         % π
\DeclareUnicodeCharacter{03C1}{\ensuremath{\rho}}        % ρ
\DeclareUnicodeCharacter{03C3}{\ensuremath{\sigma}}      % σ
\DeclareUnicodeCharacter{03C4}{\ensuremath{\tau}}        % τ
\DeclareUnicodeCharacter{03C6}{\ensuremath{\varphi}}     % φ
\DeclareUnicodeCharacter{03C7}{\ensuremath{\chi}}        % χ
\DeclareUnicodeCharacter{03C8}{\ensuremath{\psi}}        % ψ
\DeclareUnicodeCharacter{03C9}{\ensuremath{\omega}}      % ω
\DeclareUnicodeCharacter{2013}{--}              % – (en dash)
\DeclareUnicodeCharacter{2014}{---}             % — (em dash)
\DeclareUnicodeCharacter{2018}{`}               % ‘
\DeclareUnicodeCharacter{2019}{'}               % ’
\DeclareUnicodeCharacter{201C}{``}              % “
\DeclareUnicodeCharacter{201D}{''}              % ”
\DeclareUnicodeCharacter{2070}{\ensuremath{^{0}}}        % ⁰
\DeclareUnicodeCharacter{2071}{\ensuremath{^{i}}}        % ⁱ
\DeclareUnicodeCharacter{2122}{\textsuperscript{TM}}     % ™
\DeclareUnicodeCharacter{2126}{\ensuremath{\Omega}}      % Ω (ohm sign)
\DeclareUnicodeCharacter{210F}{\ensuremath{\hbar}}       % ℏ
\DeclareUnicodeCharacter{2115}{\ensuremath{\mathbb{N}}}  % ℕ
\DeclareUnicodeCharacter{2102}{\ensuremath{\mathbb{C}}}  % ℂ
\DeclareUnicodeCharacter{211D}{\ensuremath{\mathbb{R}}}  % ℝ
\DeclareUnicodeCharacter{2124}{\ensuremath{\mathbb{Z}}}  % ℤ
\DeclareUnicodeCharacter{211A}{\ensuremath{\mathbb{Q}}}  % ℚ
\DeclareUnicodeCharacter{2192}{\ensuremath{\to}}         % →
\DeclareUnicodeCharacter{21D2}{\ensuremath{\Rightarrow}} % ⇒
\DeclareUnicodeCharacter{2200}{\ensuremath{\forall}}     % ∀
\DeclareUnicodeCharacter{2203}{\ensuremath{\exists}}     % ∃
\DeclareUnicodeCharacter{2208}{\ensuremath{\in}}         % ∈
\DeclareUnicodeCharacter{2209}{\ensuremath{\notin}}      % ∉
\DeclareUnicodeCharacter{2227}{\ensuremath{\wedge}}      % ∧
\DeclareUnicodeCharacter{2228}{\ensuremath{\vee}}        % ∨
\DeclareUnicodeCharacter{2229}{\ensuremath{\cap}}        % ∩
\DeclareUnicodeCharacter{222A}{\ensuremath{\cup}}        % ∪
\DeclareUnicodeCharacter{2245}{\ensuremath{\cong}}       % ≅
\DeclareUnicodeCharacter{2248}{\ensuremath{\approx}}     % ≈
\DeclareUnicodeCharacter{2260}{\ensuremath{\neq}}        % ≠
\DeclareUnicodeCharacter{2264}{\ensuremath{\leq}}        % ≤
\DeclareUnicodeCharacter{2265}{\ensuremath{\geq}}        % ≥
\DeclareUnicodeCharacter{22C5}{\ensuremath{\cdot}}       % ⋅
\DeclareUnicodeCharacter{2295}{\ensuremath{\oplus}}      % ⊕
\DeclareUnicodeCharacter{2297}{\ensuremath{\otimes}}     % ⊗

% ---- Mathematics & symbols ------------------------------------------
\usepackage{amsmath,amssymb,bm,mathrsfs,mathtools}
\usepackage{amsthm}             % theorem environments (load BEFORE hyperref)

% ---- Theorem environments (shared counter set, TECT canonical) -------
% All theorems / lemmas / propositions / corollaries / remarks share a
% single counter so cross-references read consistently across a paper.
% Definitions get a separate counter (definitions are numbered in their
% own sequence by editorial convention).
\newtheorem{theorem}{Theorem}
\newtheorem{lemma}[theorem]{Lemma}
\newtheorem{proposition}[theorem]{Proposition}
\newtheorem{corollary}[theorem]{Corollary}

\theoremstyle{remark}
\newtheorem{remark}[theorem]{Remark}
\newtheorem{example}[theorem]{Example}

\theoremstyle{definition}
\newtheorem{definition}[theorem]{Definition}
\newtheorem{conjecture}[theorem]{Conjecture}
\newtheorem{hypothesis}[theorem]{Hypothesis}

% ---- Tables / floats / graphics --------------------------------------
\usepackage{booktabs}           % \toprule / \midrule / \bottomrule
\usepackage{array}
\usepackage{graphicx}
\usepackage{enumitem}           % \begin{itemize}[leftmargin=...] options
\usepackage{hyphenat}           % \hyp{} for breakable hyphens in \texttt
\usepackage{seqsplit}           % \seqsplit{} for long unbreakable strings
\sloppy                         % allow stretchier inter-word spacing to
                                % reduce overfull \hbox warnings on long URLs
                                % and long \texttt tokens (paragraph-level).
\emergencystretch=3em           % last-resort hbox stretch budget

% ---- Cross-references (hyperref last among the standard packages,
%      cleveref AFTER hyperref) ----------------------------------------
\usepackage[%
  colorlinks=true,%
  linkcolor=blue,%
  citecolor=blue,%
  urlcolor=blue,%
  breaklinks=true,%
  bookmarks=true,%
  bookmarksnumbered=true%
]{hyperref}
\usepackage[capitalise,nameinlink]{cleveref}

% ---- TECT-canonical author block macros -----------------------------
% (defined here so every paper has the same author / affiliation style)
\newcommand{\tectauthor}{%
  \author{Jusang Lee}%
  \email{jtkor@outlook.com}%
  \affiliation{Independent researcher, Republic of Korea}%
}

% ---- TECT-canonical archive footer ----------------------------------
\newcommand{\tectarchivefooter}{%
  \par\medskip
  \noindent\small\textit{Canonical archive}:
  \url{https://github.com/TECT-OpenPhysics/TECT}.\par%
}

% ---- TECT TODO / AUDIT-FLAG / HONEST-SCOPE inline marks --------------
% Used in drafts to highlight pending audit items without disturbing
% the body text style. Suppress via \tectsuppressmarks (define empty).
\providecommand{\tectauditflag}[1]{\textcolor{red}{\textbf{[AUDIT-FLAG: #1]}}}
\providecommand{\tecttodo}[1]{\textcolor{orange}{\textbf{[TODO: #1]}}}
\providecommand{\tecthonestscope}[1]{\textit{Honest scope: #1.}}

% =====================================================================
% End of TECT canonical preamble
% =====================================================================


\begin{document}

\title{A conditional BRST gauge-fixing and Faddeev--Popov closure note
for the TECT defect sector, with a TECT-specific Berry correction
from Math160}

\author{Jusang Lee}
\email{jtkor@outlook.com}
\affiliation{Independent researcher, Republic of Korea}

\date{\today}

\begin{abstract}
We present a candidate BRST gauge-fixing and Faddeev--Popov closure
note for the GAP-2 quantum-completion gate in the Topological Energy
Condensate Theory (TECT) defect sector. The construction combines
(a) the standard BRST / Faddeev--Popov path-integral measure for the
gauge field in covariant gauge (textbook QFT input) and (b) a
candidate TECT-specific Berry-phase correction inherited from
Math160 arising from the adiabatic evolution of the BCC order
parameter during freeze-out. The current canonical tier of GAP-2 is
\textbf{T6} per Math307 / Math310, conditional on Pillar 4's
stability theorem (Math80-AddA, hypothesis H1.1) via the
inheritance mechanism of Math280.

\textbf{Honest scope}: this draft is a conditional closure note, NOT
a stand-alone derivation. The novelty content lies in the
Berry-correction inheritance from Math160; the present draft cites
Math160 rather than re-deriving the correction, and the standard
BRST/FP textbook part is sketched at programme level rather than
fully written out with Grassmann measure and sign-convention
discipline. The over-claim wording (``GAP-2 is now complete'',
``fully rigorous'', ``ready for publication'') of an earlier draft
has been removed; promotion to a stand-alone derivation requires
(i) full BRST/FP textbook write-up with Grassmann measure, (ii)
explicit re-derivation of the Math160 Berry correction within the
present paper, (iii) consistency cross-check with the canonical
GAP-2 tier T6 in \texttt{Docs/status/TOE-FACT-SHEET.md}.
\end{abstract}

\maketitle

% =====================================================================
\section{Introduction}\label{sec:intro}
% =====================================================================
GAP-2 addresses a textbook quantum-field-theoretic issue, the
gauge-fixed path-integral measure for non-Abelian gauge theories.
The standard treatment proceeds via the Faddeev--Popov determinant
\cite{FaddeevPopov1967}, which introduces auxiliary anti-commuting
ghost fields to compensate for the over-counting introduced by
gauge fixing, and via the BRST cohomology
\cite{BRST1976,Tyutin1975}, which encodes the residual rigid
symmetry that selects the physical Hilbert subspace.

In the TECT defect-sector context, the gauge field of Pillar~4
acquires the same Faddeev--Popov / BRST structure as the SM gauge
sector; the present paper records the inheritance of that structure
under the conditional Pillar~4 closure (hypothesis H1.1 of
Math80-AddA), supplemented by a candidate TECT-specific
Berry-phase correction inherited from Math160 due to the adiabatic
evolution of the BCC order parameter through the defect manifold.
The present paper does \emph{not} re-derive the textbook BRST/FP
machinery; it inherits it under hypothesis (H1) and combines it with
the Math160 Berry inheritance under hypothesis (H2). The closure is
\emph{conditional}, not stand-alone.

% =====================================================================
\section{Setting --- BRST / Faddeev--Popov inheritance in covariant
gauge}\label{sec:setting}
% =====================================================================

We work throughout in the Minkowski signature $(+,-,-,-)$ with
metric $\eta_{\mu\nu}$ and adopt the standard non-Abelian
Yang--Mills action
\begin{equation}\label{eq:SYM}
  S_{\mathrm{YM}}[A] \;=\;
  -\frac{1}{4}\int d^4x\, F^{a}_{\mu\nu}\,F^{a\,\mu\nu},
  \qquad
  F^{a}_{\mu\nu} = \partial_\mu A^{a}_\nu - \partial_\nu A^{a}_\mu
                 + g\,f^{abc} A^{b}_\mu A^{c}_\nu,
\end{equation}
with $\{T^a\}$ the gauge-group generators in the matter
representation, $f^{abc}$ the structure constants, and $a,b,c$
running over the adjoint index of the Pillar~4 gauge sector.

The covariant gauge-fixing functional is
$G^{a}(A) = \partial^{\mu}A^{a}_{\mu}$, with parameter $\xi$. The
Faddeev--Popov \cite{FaddeevPopov1967} insertion of the identity
\begin{equation}\label{eq:FP-identity}
  1 \;=\;
  \det\!\Bigl(\tfrac{\delta G^{a}(A^{\omega})}{\delta \omega^{b}}\Bigr)
  \,\int \mathcal{D}\omega\ \delta\bigl(G^{a}(A^{\omega}) - \chi^{a}\bigr),
\end{equation}
followed by Gaussian-averaging $\chi^{a}$ with weight
$\exp\!\bigl[-\tfrac{i}{2\xi}\!\int d^4x (\chi^a)^2\bigr]$ and
exponentiation of the determinant via auxiliary anti-commuting
(Grassmann) ghost fields $c^{a},\bar c^{a}$, gives the standard
gauge-fixed gauge generating functional
\begin{equation}\label{eq:Z-gauge-fixed}
  Z[J,\bar\eta,\eta]
  \;=\;
  \int \mathcal{D}A\,\mathcal{D}\bar c\,\mathcal{D}c\;
  \exp\!\Bigl(
    i\,S_{\mathrm{eff}}[A,\bar c,c]
    + i\!\int d^4x\bigl(J^{a\,\mu}A^{a}_{\mu}
       + \bar\eta^{a} c^{a} + \bar c^{a}\eta^{a}\bigr)
  \Bigr),
\end{equation}
with the gauge-fixed effective action
\begin{equation}\label{eq:Seff}
  S_{\mathrm{eff}}[A,\bar c,c]
  \;=\;
  S_{\mathrm{YM}}[A]
  \;-\;\frac{1}{2\xi}\!\int d^4x\,\bigl(\partial^{\mu}A^{a}_{\mu}\bigr)^{2}
  \;-\;\!\int d^4x\,\bar c^{a}\,\partial^{\mu}D_{\mu}^{ac}\,c^{c},
\end{equation}
where $D_{\mu}^{ac} = \partial_{\mu}\delta^{ac} + g\,f^{abc}A^{b}_{\mu}$
is the adjoint covariant derivative. The ghost integration measure
$\mathcal{D}\bar c\,\mathcal{D}c$ is the textbook Berezin (Grassmann)
measure, with the convention
$\int d\bar c\,dc\,\bar c\,c = -1$ per pair so that the
exponentiation \eqref{eq:Z-gauge-fixed} reproduces the determinant in
\eqref{eq:FP-identity} with a positive sign.

\begin{theorem}[BRST / FP closure of the Pillar~4 gauge sector with
TECT-specific Math160 Berry correction;
\textbf{T6 \textsc{proved conditional} on (H1)+(H2)}]\label{thm:main}
Under the hypotheses
\begin{itemize}
  \item[(\textbf{H1})] the Pillar~4 gauge-stability theorem of
    Math80-AddA (hypothesis H1.1, the Grassmannian-stability
    component of the Math80 closure programme);
  \item[(\textbf{H2})] the Math160 Berry-correction inheritance
    theorem (the BCC order parameter undergoes adiabatic evolution
    through the Pillar~4 defect manifold and contributes a
    geometric phase to the ghost effective action),
\end{itemize}
the gauge-fixed effective action of \eqref{eq:Seff} acquires the
TECT-specific Berry contribution and reads
\begin{equation}\label{eq:Seff-TECT}
  \boxed{\;
  S_{\mathrm{eff}}^{\mathrm{TECT}}[A,\bar c,c]
  \;=\;
  S_{\mathrm{YM}}[A]
  \;-\;\frac{1}{2\xi}\!\int d^4x\,(\partial^{\mu}A^{a}_{\mu})^{2}
  \;-\;\!\int d^4x\,\bar c^{a}\,\partial^{\mu}D_{\mu}^{ac}\,c^{c}
  \;+\;
  S^{\mathrm{Berry}}_{\mathrm{TECT}}[A,\bar c,c],
  \;}
\end{equation}
where $S^{\mathrm{Berry}}_{\mathrm{TECT}}$ is the Math160-derived
Berry-phase correction. Under (H1)+(H2), the action
\eqref{eq:Seff-TECT} is invariant under the standard nilpotent BRST
transformation
\begin{equation}\label{eq:BRST}
  s A^{a}_\mu = D^{ac}_\mu c^{c}, \quad
  s c^{a} = -\tfrac{1}{2}g\,f^{abc}c^{b}c^{c}, \quad
  s \bar c^{a} = -\tfrac{1}{\xi}\,\partial^{\mu}A^{a}_{\mu},
  \quad s^{2}=0,
\end{equation}
so that the BRST cohomology of \eqref{eq:Seff-TECT} agrees with that
of the textbook gauge-fixed Yang--Mills action; in particular, the
physical-state subspace is the BRST-closed-modulo-exact subspace of
the total Hilbert space, and the $S$-matrix on physical states is
unitary.

\textbf{Conditional tier:} GAP-2 is recorded at \textbf{T6
\textsc{proved conditional}} per Math307 / Math310, conditional on
(H1) Pillar~4 H1.1 \emph{and} (H2) the Math160 Berry inheritance.
\end{theorem}

\begin{remark}[Honest scope of the theorem]\label{rmk:honest}
Theorem~\ref{thm:main} \emph{inherits} the textbook BRST / FP
machinery rather than re-deriving it within the present paper, and
\emph{cites} the Math160 Berry-correction derivation rather than
re-deriving it. The novelty content of the present paper is the
inheritance argument under (H1)+(H2), not the textbook-side
construction. Promotion of GAP-2 to a stand-alone derivation
requires (i) full BRST/FP textbook write-up with explicit Grassmann
measure, sign, and Wick-rotation conventions; and (ii) explicit
re-derivation of $S^{\mathrm{Berry}}_{\mathrm{TECT}}$ within the
present paper. Neither (i) nor (ii) is performed here; the
conditional T6 verdict reflects the inheritance, not a stand-alone
derivation.
\end{remark}

% =====================================================================
\section{Proof sketch under (H1)+(H2)}\label{sec:proof}
% =====================================================================

\paragraph{Step 1 (BRST symmetry setup).}
The nilpotent BRST transformation in the conventions of
\eqref{eq:BRST} is the textbook construction
\cite{BRST1976,Tyutin1975}. It satisfies $s^2 = 0$ on every field
(this is the Jacobi identity for $f^{abc}$ together with the
Grassmann anti-commutativity of $c^{a}$).

\paragraph{Step 2 (Faddeev--Popov implementation).}
Insertion of the identity \eqref{eq:FP-identity} into the
unrestricted gauge-field path integral, Gaussian-averaging the
gauge-fixing functional, and exponentiating the determinant via the
Berezin (Grassmann) measure for $c^{a},\bar c^{a}$ yields the
gauge-fixed effective action \eqref{eq:Seff}.

\paragraph{Step 3 (Berry-phase inheritance from Math160).}
Math160 derives the TECT-specific Berry-phase contribution
$S^{\mathrm{Berry}}_{\mathrm{TECT}}[A,\bar c, c]$. The contribution
arises because the BCC order parameter undergoes adiabatic evolution
through the Pillar~4 defect manifold during phase-transition
freeze-out and acquires a geometric phase encoded in the ghost
boundary conditions on the defect manifold. Under hypothesis (H2)
the resulting contribution is added to \eqref{eq:Seff} producing
\eqref{eq:Seff-TECT}.

\paragraph{Step 4 (BRST closure verification).}
Under hypothesis (H1), $S^{\mathrm{Berry}}_{\mathrm{TECT}}$ is
$s$-closed (a property derived in Math160 from the gauge-equivariance
of the underlying Berry connection). Therefore the total action
\eqref{eq:Seff-TECT} is BRST-invariant, $s\,S^{\mathrm{TECT}}_{\mathrm{eff}}=0$,
and the BRST cohomology of \eqref{eq:Seff-TECT} agrees with that of
the textbook gauge-fixed Yang--Mills action.

\paragraph{Step 5 (Conditional inheritance from Pillar~4 via
Math280).}
Math280 establishes the inheritance mechanism: under (H1) Pillar~4
hypothesis H1.1, the BRST/FP closure of \eqref{eq:Seff-TECT}
inherits the Pillar~4 tier. With Pillar~4 at the present canonical
T6 \textsc{proved conditional}, the GAP-2 closure recorded here is
\textbf{T6 \textsc{proved conditional} on (H1)+(H2)}.
No additional hypothesis beyond (H1)+(H2) is required.

% =====================================================================
\section{Numerical estimate of the Berry correction}
\label{sec:numerics}
% =====================================================================

In the covariant gauge with fixing parameter $\xi = 1$ (Feynman
gauge) and at the linearised level $g\to 0$, the ghost piece of
\eqref{eq:Seff-TECT} reduces to
\begin{equation}\label{eq:Lghost-feynman}
  \mathcal{L}_{\mathrm{ghost}}^{(\xi=1,\,g=0)}
  \;=\;
  -\,\bar c^{a}\,\partial^{\mu}\partial_{\mu}\,c^{a}
  \;+\;\delta\mathcal{L}^{\mathrm{Berry}}_{\mathrm{TECT}},
\end{equation}
where $\delta\mathcal{L}^{\mathrm{Berry}}_{\mathrm{TECT}}$ is the
linearised limit of the Math160 Berry contribution. Numerical
evaluation in Math160 gives a relative correction of order
$\sim 10^{-2}$ to the ghost propagator at the one-loop matching
scale; this is consistent with the $\mathcal{O}(a_{\mathrm{BCC}}^{2})$
expectation for a continuum-limit-vanishing geometric correction.
The standard Faddeev--Popov ghost propagator is recovered to leading
order in $a_{\mathrm{BCC}}$.

% =====================================================================
\section{Discussion --- conditional tier, not stand-alone derivation}
\label{sec:discussion}
% =====================================================================

\paragraph{What the present note records.}
Under hypotheses (H1)+(H2), the gauge-fixed action of the Pillar~4
gauge sector takes the form \eqref{eq:Seff-TECT} and admits a
nilpotent BRST symmetry \eqref{eq:BRST}, so the BRST/FP cohomology
agrees with that of the textbook gauge-fixed Yang--Mills sector. The
GAP-2 canonical tier per Math307 / Math310 is \textbf{T6 \textsc{proved
conditional}} on (H1) Pillar~4 H1.1 \emph{and} (H2) the Math160
Berry-correction inheritance. Neither hypothesis is discharged in the
present paper.

\paragraph{What the present note does NOT establish.}
\begin{enumerate}
  \item It does \emph{not} establish a stand-alone BRST/FP derivation
    of GAP-2: the textbook part is inherited, the Math160 Berry part
    is inherited (Remark~\ref{rmk:honest}).
  \item It does \emph{not} discharge Pillar~4 hypothesis H1.1: that
    is the canonical Pillar~4 sub-task tracked separately in
    Math80-AddA.
  \item It does \emph{not} certify the asymptotic behaviour of
    $S^{\mathrm{Berry}}_{\mathrm{TECT}}$ in the continuum limit
    (lattice spacing $\to 0$); the Math160 Berry correction is
    expected to vanish in that limit at order
    $\mathcal{O}(a_{\mathrm{BCC}}^2)$, but a continuum-limit
    constant-bound theorem on this is a separate item.
\end{enumerate}

\paragraph{Promotion path to stand-alone derivation.}
A stand-alone derivation of GAP-2 within a future revision of this
paper requires (i) explicit textbook write-up of the BRST/FP machinery
including the Grassmann measure, the Euclidean--Minkowski Wick
rotation, and the $i\epsilon$ prescription; (ii) explicit
re-derivation of $S^{\mathrm{Berry}}_{\mathrm{TECT}}$ within the
present paper; (iii) cross-check with the canonical GAP-2 tier T6 in
\texttt{Docs/status/TOE-FACT-SHEET.md} after promotion. Until those
items are completed, the present note remains a \emph{conditional
closure note}, not a stand-alone derivation.

% =====================================================================
\begin{acknowledgments}
The author thanks Math160 for the BRST formulation and Berry-phase
derivation. This work is part of TECT, canonical archive at
\url{https://github.com/TECT-OpenPhysics/TECT}.
\end{acknowledgments}

% =====================================================================
\bibliographystyle{apsrev4-2}
\bibliography{../_shared/tect-references}

\end{document}
